%===============================================================================
%
%===============================================================================
\section{Écrire du code pour ManuX}

   Le but de cette section est de montrer brièvement comment écrire du
code dans le noyau \manux.

   Je vais partir du principe que vous souhaitez ajouter une nouvelle
fonctionalité à ManuX (c'est vexant, vous sous-entendez que ManuX n'a
pas tout !?). Dans ce cas-là, le mieux est de construire une nouvelle
version du noyau minimaliste puis d'y ajouter ce dont vous avez
besoin.

%-------------------------------------------------------------------------------
%
%-------------------------------------------------------------------------------
\subsection{Les bases pour construire un noyau vide}

   Imaginons que vous souhaitiez ajouter la fonctionalité `toto` au
noyau. Vous allez donc chercher à construire un noyau que vous
nommerez par exemple `test-toto`. Il faudra pour cela créer deux
fichiers :

\begin{itemize}
   \item Le fichier `multiconf/test-toto.h` décrit la configuration du
     noyau. Vous pourrez par exemple partir de `multiconf/nop.h` que
     vous enrichirez ensuite.
   \item Le fichier `noyau/main-test-toto.c` est quant à lui le
     fichier d'initialisation du noyau. Vous pouvez là aussi partir de
     `noyau/main-nop.c`. Le nom de ce fichier est défini dans
     `multiconf/test-toto.h`.
\end{itemize}

   Suite à la création de ces deux fichiers, vous pourrez alors
utiliser `make` pour compiler et lancer votre noyau :

\begin{lstlisting}
  $ make CFG=test-toto
  $ make CFG=test-toto run
\end{lstlisting}
